% 由已核对的中文初译逐句精修；完整英文底稿与审计另存于 paper/closeout。
\section{恒定最长弦与成对径向剖面}\label{sec:constant-length-classification}

% [CL01]
第~\ref{sec:attainment} 节留下了一个分类问题：如果最长完整弦在一个方向区间上始终没有长度余量，体的径向剖面会受到什么限制？一条弦可能跨过几个生成片，其内部的角点也可能阻止载线移动。下面利用实际边界的有限结构，逐一处理端点与内部角点的接触。

% [CL02]
设 $K\subset\mathbb R^2$ 是紧的半代数集，包含原点，并且关于原点星形，即每个 $x\in K$ 都满足 $[0,x]\subset K$。假设它的实际边界可以有限分割成相对开直线段、以原点为圆心且半径为正的相对开圆弧，以及点。在交点和端点处切分这些片；同一直线或圆上的重叠片先合并，再切分。因此，除去有限个顶点，每个边界点都有一个邻域，其中边界只是一条光滑直线或圆弧。这允许径向尖刺，也允许原点在边界上；不要求 $K$ 凸，也不要求它包含原点的某个邻域。

% [CL03]
对射影方向 $q\in\mathbb {RP}^1$，在局部选取单位提升 $u=u_q$，并令 $n=n_q$ 为它逆时针旋转直角所得的垂直向量。定义
\[
 L_K(q)=\max\{\ell\ge0:[x,x+\ell u]\subset K\text{，对某个 }x\in K\},
 \qquad \rho_K(u)=\max\{t\ge0:tu\in K\}.
\]
定义 $L_K$ 的最大值针对完整的连通线段，而不是 $K$ 中任意两点之间的距离。紧致性保证两个最大值都存在。星形性给出完整的中心截面
\[
 K\cap\mathbb Ru=[-s(q)u,r(q)u],\qquad
 r(q)=\rho_K(u),\quad s(q)=\rho_K(-u),\quad r(q)+s(q)\le L_K(q).
\]
更换提升会交换 $r$ 和 $s$；它们的和以及成对平方和不依赖于提升。

% [CL04]
\begin{lemma}[恒长区间上的分类]\label{lem:constant-length-classification}
对满足上述假设的体，设 $L_K(q)=1$ 在非空开射影区间 $J$ 上成立。那么每个 $q\in J$ 都有一条位于过原点直线上的最长单位段；等价地，$r(q)+s(q)=1$ 在 $J$ 上成立。存在有限的\emph{方向}集 $E\subset J$，使得在 $J\setminus E$ 的每个连通分支上，对连续选取的提升，$r$ 和 $s$ 分别是常数 $r_I,s_I\ge0$，满足 $r_I+s_I=1$。每个正的常数都来自一条以原点为圆心的实际边界圆弧。结论允许其中一个常数为零。
\end{lemma}

\begin{proof}
% [CL05]
先把最长段归入有限种接触类型。以 $V$ 表示边界分割的所有顶点，并把原点加入 $V$。将平行于某个直边的方向，以及连接 $V$ 中两个不同点的方向，全部放入例外集。这两类都是有限集。在下面的论证中，我们还可以加入有限个细分方向；绝不删去一个方向区间。

% [CL06]
在半径为 $R$ 的光滑圆弧边界点附近，$K$ 局部位于圆盘内侧。事实上，邻近圆弧属于闭集 $K$，所以它的径向填充包含一个向内的开扇形。在足够小的邻域中，这段圆弧是唯一的边界；在圆弧两侧各自连通的区域内，是否属于 $K$ 保持不变。内侧属于 $K$，外侧不可能也属于 $K$，否则圆弧上的点就是内点。因此，一条包含于 $K$ 的非平凡线段不能以内点接触这段圆弧：横截会进入外侧，而切线在接点两侧都位于圆盘外。同样的切线理由表明，线段若以圆弧上的点为端点，也必须横截圆弧。在光滑直边点处，线段的内部接触只能沿着边界直线发生，除非该点是 $K$ 的内点；后者与该点属于边界矛盾。平行方向已经去掉。如果直边两侧都没有 $K$ 的内点，例如在尖刺处，横截它的线段也不可能包含于 $K$。

% [CL07]
最长正长度线段的两个端点都在边界上，否则可以从内点端继续延长。避开上述有限方向后，只剩两种可能。它的载线经过某个 $v\in V$，这个点可以是端点，也可以是内部接触点；载线上至多有一个这样的点。这是固定顶点情形。否则，整个开线段都在 $\operatorname{int}K$ 中，两个端点都是光滑边界片上的横截点。这是自由平移情形。后一种载线作小幅平移后，仍有一条包含线段，端点落在原来的两个边界片上：紧的中间子段到补集有正距离，而横截性和局部内侧保证端点附近仍然包含。因此，这条邻近线段的长度关于载线的带符号高度可微，并在最长段处驻定。这里的论证针对 $K$ 中的包含关系，而不是某一个生成片中的包含关系。

% [CL08]
接触类型还须在 $J$ 上作有限细分后固定下来。关系 $[x,y]\subset K$ 是半代数的：它可写成对所有 $t\in[0,1]$ 都有 $x+t(y-x)\in K$。量词消去也使最长段关系成为半代数关系。半代数选择在每个局部方向图上提供一条有序的最长段。对它的图像和接触关系作有限胞腔分解，便可在有限个方向处细分 $J$，使每个剩余开区间上所选线段都具有固定的接触类型：相同的端点边界片、相同的可能顶点，以及相同的圆交点符号。单变量半代数函数可有限分割为连续光滑片；它的零集同样由有限个点和区间组成。我们在这些片的端点处分割。在整个区间上持续存在的退化必须另作处理，不能当作例外删掉。过固定非零顶点 $v$ 的直线若在一个区间上始终与半径为 $R$ 的圆相切，就要求 $(v\cdot u)^2=|v|^2-R^2$ 在整个区间上成立，这是不可能的。孤立的相切方向可以去掉。中心直线与每个正半径圆都是横截的。端点落在顶点的情形会在下面明确保留。

% [CL09]
先考虑一个自由平移区间。把载线写成 $x=tu+hn$。圆端点和直边端点的纵向坐标分别为
\[
 t_C=\sigma\sqrt{R^2-h^2},\qquad
 t_D=\frac{c-hB}{A},\qquad
 A=a\cdot u,\quad B=a\cdot n,\quad |a|=1,
\]
其中直边的载线为 $a\cdot x=c$，而 $\sigma\in\{-1,1\}$。横截性保证 $R^2-h^2>0$ 且 $A\ne0$。有序弦长为 $t_2-t_1=1$，其高度导数为零。

% [CL10]
对两个圆端点，这个导数方程为
\[
 h\left(\frac{\sigma_1}{\sqrt{R_1^2-h^2}}
          -\frac{\sigma_2}{\sqrt{R_2^2-h^2}}\right)=0.
\]
若 $h\ne0$，两个分式相等就迫使符号相同、半径相等，从而纵向坐标相同，弦长为零。因此 $h=0$，最长段位于中心直线上。$K$ 的中心截面包含它，并且长度至多为一，所以这个截面本身长一且包含原点。

% [CL11]
对两个直边端点，驻值条件给出 $B_1/A_1=B_2/A_2$。两个法向量因此平行。调整它们的符号，使它们成为同一个单位向量 $a$。两个端点的带符号差便是 $(c_2-c_1)/(a\cdot u)$。非零弦满足 $c_2-c_1\ne0$；这个表达式不可能在开方向区间上恒等于一，因为非零线性泛函 $a\cdot u$ 在这样的区间上不为常数。因此，该接触类型不能在所选恒长区间上持续存在。

% [CL12]
对一个圆端点和一个直边端点，驻值条件使 $-h/t_C$ 与 $-B/A$ 相等，所以 $hA=t_CB$。因此圆端点 $t_Cu+hn$ 平行于 $a=Au+Bn$，并等于 $\tau Ra$；细分后 $\tau\in\{-1,1\}$ 固定。代入得到
\[
 t_C-t_D=\frac{\tau R-c}{a\cdot u}.
\]
根据端点次序，这个值必须为一或负一。分子为零时长度为零；分子非零时，又会使 $a\cdot u$ 为常数。这排除了持续存在的混合自由平移区间，也包括两个端点属于不同生成片的弦。

% [CL13]
剩下的是固定顶点区间。把载线写成 $x=v+tu$。若 $v=0$，载线已经是中心直线，可以应用前面的中心截面论证。以下假设 $v\ne0$。圆端点与直边端点的坐标为
\[
 t_C=-v\cdot u+\sigma\sqrt{R^2-|v|^2+(v\cdot u)^2},
 \qquad t_D=\frac{d}{T},\qquad d=c-a\cdot v,\quad T=a\cdot u.
\]
经过上述相切方向细分，平方根严格为正，且 $T\ne0$。即使 $v$ 是线段内部的凹角接触点，两个端点都不等于 $v$，这些公式仍然描述了端点。

% [CL14]
对两个圆端点，置 $z=(v\cdot u)^2$、$C_i=R_i^2-|v|^2$。它们的差为
\[
 \sigma_2\sqrt{C_2+z}-\sigma_1\sqrt{C_1+z}.
\]
当 $v\ne0$ 时，函数 $z$ 在任何开方向区间上都不是常数；缩小区间后，可以用它作参数。若端点差为常数，它关于 $z$ 的导数就必须为零。符号相反时导数不可能为零。符号相同时，导数为零迫使 $C_1=C_2$，于是端点差为零，而不是一。因此这种情形不可能出现。

% [CL15]
对两个直边端点，置 $d_i=c_i-a_i\cdot v$、$T_i=a_i\cdot u$。长度方程给出
\[
 d_2T_1-d_1T_2=T_1T_2.
\]
三角多项式恒等式若在一个开弧上成立，就在整个单位圆上成立。作 $u\mapsto-u$ 后，左边是奇的，右边是偶的，所以该恒等式将迫使 $T_1T_2$ 处处为零。两个非零线性泛函不可能具有这一性质。这就排除了直边与直边的情形，也包括 $d_i=0$ 的情况。

% [CL16]
对一个圆端点和一个直边端点，它们的带符号差是固定的 $\varepsilon\in\{-1,1\}$。因此圆端点为 $v+(d/T+\varepsilon)u$。把它的圆方程乘以 $T^2$，得到
\begin{equation}\label{eq:classification-mixed-identity}
 (|v|^2-R^2)T^2+(d+\varepsilon T)^2
       +2(v\cdot u)(d+\varepsilon T)T=0.
\end{equation}
这同样是整个圆上的三角多项式恒等式。在垂直于 $a$ 的单位向量处代入，得到 $d^2=0$。在原来的区间上，$T\ne0$，恒等式于是化为 $|v+\varepsilon u|^2=R^2$。这表示 $v\cdot u$ 在开弧上为常数，迫使 $v=0$，与当前情形矛盾。被延拓到 $T=0$ 的只有多项式恒等式；没有假设几何端点分支在那里也存在。

% [CL17]
若一个端点就是固定顶点本身，它的坐标为 $t=0$，另一个端点的坐标必须是 $\varepsilon\in\{-1,1\}$。另一个端点在圆上时，得到 $|v+\varepsilon u|^2=R^2$，再次迫使 $v=0$。另一个端点在直边上时，得到 $d=\varepsilon(a\cdot u)$，不可能在开弧上成立。若两个端点都是顶点，对应方向已在最初去掉；被两个不同内部顶点约束的载线也已同时排除。这些情形包含了所有固定端点和内部固定顶点的构型。

% [CL18]
我们已经证明，在每个剩余开区间上，所选最长段的载线都经过原点。因此，它的完整中心截面长一。对例外方向 $q\in J$，取非例外方向 $q_j\to q$ 及相应的长一中心截面，并选择相容的提升。由于 $K$ 有界，两个非负半径有一个收敛子序列；由闭性，极限线段仍包含在 $K$ 中，长一且包含原点。由于 $L_K(q)=1$，方向 $q$ 的中心截面不可能更长。因此 $r(q)+s(q)=1$ 在 $J$ 的每个方向都成立，包括所有例外方向。

% [CL19]
最后再细分方向区间，确定两个径向半径的形式。正的径向端点属于边界，否则可以沿射线继续延伸。除去非零顶点对应的有限方向和有限个半代数细分点后，这个端点落在一个固定的光滑边界片上。它的半径或者是圆弧常数 $R>0$，或者是直边表达式 $c/(a\cdot u)$；与该直边平行的射线已经排除。半径为零也作为一种可能在整个区间上存在的类型保留下来。将 $u$ 换成 $-u$，就可同样处理负端点。

% [CL20]
若一个半径为常数，恒等式 $r+s=1$ 就使另一个半径也为常数。正的直边表达式不可能在开区间上为常数：除非 $c=0$，否则正的常数会要求 $a\cdot u$ 为常数；而 $c=0$ 给出的是零半径，不是正的径向端点。若两个正端点都在直边上，它们的和将具有形式
\[
 \frac{c_1}{a_1\cdot u}-\frac{c_2}{a_2\cdot u}=1.
\]
清去分母后，关于 $u$ 的奇一次式将等于偶二次式。直边与直边情形中的论证已经排除了这一可能。若一个半径为零，另一个就等于一，因此来自圆弧。于是每个剩余区间上的两个半径都是常数，且每个正半径都由圆弧承载。引理得证。
\end{proof}

% [CL21]
\subsection*{为什么原始有限体具有这种边界}
现在检查原始有限体确实满足边界假设。对~\eqref{eq:upper-original} 中的每条记录分别考察。一条记录把固定三角形 $T=\operatorname{conv}(0,A,B)$ 在闭角区间 $I$ 上旋转。它在物理角 $\phi$ 处的径向剖面，是 $q\in I$ 上 $\rho_T(\phi-q)$ 的最大值。对非退化三角形，在相邻顶点方向之间，$\rho_T(\theta)$ 的每个正光滑分支都具有直边形式 $c/(a\cdot u_\theta)$。正分支的内部驻点具有正的二阶导数，所以这样的点是严格局部极小点，而不是极大点。因此，关于 $q$ 的最大值必在 $I$ 的端点或 $T$ 的某个顶点方向取得。第一种情形描出端时刻旋转三角形的边，第二种情形描出以原点为圆心的顶点圆。剖面换分支、跳到零或结束的方向只有有限个，其径向边界部分是直线段。退化三角形由一条或两条径向线段组成；其扫掠是圆扇形及径向侧边，具有相同的边界类型。单点区间只给出未经扫掠的三角形或线段。

% [CL22]
在单位圆上表示方向时，这些最大值及其分支比较都是半代数的；有限条角区间弧本身也是半代数的。因此，上述细分是有限的。有限条记录的并仍然紧、半代数且关于原点星形。其边界包含在各条记录的边界之并中：在这个边界之并外，一个点或者属于某条记录的内部，或者在所有记录之外。把有限条直线和圆在交点以及暴露状态改变处分割，就恰好得到上面要求的实际边界分割。隐藏的生成圆弧不为这个论证提供边界接触。

% [CL23]
\begin{corollary}[全圆紧约束情形]\label{cor:classification-full-circle}
设 $K$ 满足引理~\ref{lem:constant-length-classification} 的边界假设，并且在每个射影方向上都包含完整单位段。定义
\[
 L_*(q)=\lim_{\delta\downarrow0}\inf_{\operatorname{dist}(p,q)<\delta}L_K(p),
 \qquad Z=\{q:L_*(q)=1\}.
\]
若 $Z=\mathbb {RP}^1$，则除有限个方向外有 $L_K=1$，并且
\begin{equation}\label{eq:classification-area}
 |K|=\frac12\int_0^\pi\bigl(r(q)^2+s(q)^2\bigr)\,dq
       \ge\frac\pi4.
\end{equation}
特别地，这样的原始有限体不能达到 $A_*$。
\end{corollary}

\begin{proof}
% [CL24]
端点的紧致性表明 $L_K$ 上半连续：对一个收敛的方向序列及其最长段，可以选出端点收敛的子序列，极限是极限方向上包含于 $K$ 的线段。它未必下半连续，因为孤立的较长弦可能在旋转后消失。引理中使用的完整线段公式说明 $L_K$ 是半代数的。它的包络也具有半代数性。为此，可在射影圆周上使用一个相容的半代数度量，包络不会改变：条件 $L_*(q)>a$ 等价于存在 $b>a$ 和 $\delta>0$，使得只要 $\operatorname{dist}(p,q)<\delta$ 就有 $L_K(p)>b$。对这个公式可以应用量词消去。此外，$L_*$ 下半连续且至少为一，所以 $Z$ 闭且半代数，从而是有限个闭区间与点的并。另一方面，$L_K$ 在紧射影圆周上只有有限个不连续方向。在每个连续方向处，它的下半连续包络都等于它自身的值。由于 $L_K\ge1$ 且处处 $L_*=1$，可知每个开连续区间上都有 $L_K=1$。现在应用引理~\ref{lem:constant-length-classification}，得到除有限个方向外 $r+s=1$。半代数性由此把包络的等式转为实际最长弦的等式，但有限个不连续方向仍须保留。

% [CL25]
紧星体的径向剖面有界且可测。极坐标积分在角长度为 $\pi$ 的射影区间上配对两个物理提升，得到~\eqref{eq:classification-area} 中的等式。在几乎每个方向上，引理给出
\[
 r^2+s^2=\frac{(r+s)^2+(r-s)^2}{2}\ge\frac12.
\]
这就证明了下界。尚未得到关系 $r+s=1$ 的有限个方向，在这个积分中只对应有限条物理射线。例外方向上的一条较长弦可能生成正面积三角形；该三角形仍属于 $K$，并已计入它的实际径向剖面。这里没有因为弦方向是例外方向就删去任何三角形。这个论证也没有把 $K$ 认定为圆盘。

% [CL26]
不达到性还需要一个更小的合法竞争体。取以原点为中心、边长为 $2/\sqrt3$ 的等边三角形 $T$，便可显式完成这个比较。它的差体 $T-T$ 是正六边形，六个顶点是三角形不同顶点之差。其外接圆半径为 $2/\sqrt3$，内切圆半径为一，所以包含闭单位圆盘。因此，每个单位向量都是 $T$ 中两点之差；凸性保证连接它们的完整单位段在 $T$ 中。它的面积为 $1/\sqrt3$，严格小于 $\pi/4$。严格不等式可直接用面积比较说明：内接于单位圆盘的正六边形面积为 $3\sqrt3/2$，严格小于 $\pi$，而 $3\sqrt3/2>4/\sqrt3$。调用第~\ref{sec:extremal} 节的有限开包络恢复，把误差取为小于 $\pi/4-1/\sqrt3$，便得到面积小于 $\pi/4$ 的原始有限体。因此，全圆紧约束情形中的原始有限体 $K$ 不可能实现原始有限类的下确界。
\end{proof}

% [CL27]
这个推论只适用于 $Z=\mathbb {RP}^1$。对混合的 $Z$，同样的半代数论证可在它的开区间片上给出常半径的中心截面，只需去掉有限个方向。这些截面的连续并由扇区组成，可以有正面积。第~\ref{sec:attainment} 节使用的闭选择器图像还必须保留例外方向及所有单侧极限单位段，包括从 $Z$ 外趋近的极限；它们可以贡献非退化的原点三角形。这些扇区与有限个例外三角形可能已经填满可用的核心。分类既没有证明混合情形中的严格核心面积差，也没有证明有限体一概不能达到下确界。混合核心刚性、全局有限不达到性以及 $A_*$ 的值仍然开放。
